\section{Detalle del proceso de difusión hacia adelante}

\subsection{Definición de la distribución de transición hacia adelante}

Cada paso de transición del proceso de difusión se define como una distribución gaussiana de forma específica:

$$q(x_t | x_{t-1}) = \mathcal{N}\big(x_t;\, \sqrt{1 - \beta_t}\, x_{t-1},\; \beta_t I\big)$$

Los significados de los diversos símbolos son los siguientes:
\begin{itemize}
    \item $\beta_t \in (0, 1)$: parámetro del cronograma de ruido (noise schedule) para el paso $t$, que es una función \textbf{no decreciente} con respecto al tiempo $t$.
    \item $\sqrt{1 - \beta_t}\, x_{t-1}$: media — escalar el valor del paso anterior hacia el origen.
    \item $\beta_t I$: varianza — inyectar ruido gaussiano isótropo.
\end{itemize}

\begin{importantbox}{Comprensión intuitiva del proceso de difusión hacia adelante}
Cada transición hacia adelante en el proceso realiza dos cosas: (1) multiplicar el valor actual $x_{t-1}$ por un coeficiente ligeramente menor que 1, $\sqrt{1-\beta_t}$, para contraerlo hacia el origen; y (2) agregar ruido gaussiano con varianza $\beta_t$. Cuando $\beta_t$ se acerca a 0, $x_t \approx x_{t-1}$ (cambia casi nada); cuando $\beta_t$ se acerca a 1, $x_t$ se vuelve casi puro ruido. A medida que $t$ aumenta, $\beta_t$ crece gradualmente y el ruido se acumula progresivamente, hasta que $x_T$ tiende a una distribución gaussiana estándar.
\end{importantbox}

\subsection{Distribución de salto hacia adelante (Forward Jump Distribution)}

Una ventaja clave del proceso de difusión hacia adelante es que no necesitamos ejecutarlo paso a paso; podemos muestrear directamente $x_t$ en cualquier paso temporal $t$ partiendo de $x_0$. Definamos $\alpha_t = 1 - \beta_t$ y $\bar{\alpha}_t = \prod_{s=1}^{t} \alpha_s$, entonces:

$$q(x_t | x_0) = \mathcal{N}\big(x_t;\, \sqrt{\bar{\alpha}_t}\, x_0,\; (1 - \bar{\alpha}_t) I\big)$$

Explicación de los símbolos:
\begin{itemize}
    \item $\bar{\alpha}_t = \prod_{s=1}^{t}(1-\beta_s)$: coeficiente acumulado conservado, que disminuye con $t$.
    \item $\sqrt{\bar{\alpha}_t}\, x_0$: media — versión escalada de los datos originales.
    \item $(1 - \bar{\alpha}_t) I$: varianza — cantidad total de ruido acumulado.
\end{itemize}

\begin{importantbox}{Importancia de la distribución de salto hacia adelante}
La distribución de salto hacia adelante $q(x_t|x_0)$ permite que, durante el entrenamiento, \textbf{directamente} muestremos $x_t$ en cualquier paso temporal a partir de $x_0$, sin necesidad de ejecutar el proceso de difusión paso a paso. Este es el factor clave para la eficiencia del entrenamiento de DDPM. Usando técnicas de reparametrización:
$$x_t = \sqrt{\bar{\alpha}_t}\, x_0 + \sqrt{1 - \bar{\alpha}_t}\, \epsilon, \quad \epsilon \sim \mathcal{N}(0, I)$$
\end{importantbox}

Esta fórmula establece la relación exacta entre $x_0$, $x_t$ y el ruido $\epsilon$, y se volverá a utilizar repetidamente al derivar el objetivo de predicción de ruido.

\subsection{Resumen de la sección}

El proceso de difusión inyecta ruido en los datos gradualmente mediante distribuciones gaussianas de transición predefinidas. Dado que la superposición de distribuciones gaussianas sigue siendo gaussiana, podemos deducir una distribución de salto hacia adelante cerrada $q(x_t|x_0)$, lo que nos permite muestrear $x_t$ directamente desde $x_0$ en un solo paso durante el entrenamiento, sin necesidad de iterar.

%% ========== Sección 3: Descomposición del ELBO y objetivo de entrenamiento ==========


